%%%%%%%%%%%%%%%%%%%%%%%%%%%%%%%%%%%%%%%%%%%%%%%%%%%%%%%%%%%%%%%%%%%%%%%%%%%%%
%% Chapter 1 — Introduction
%% Budget: 3 minutes / 3 slides
%%   S1  Motivation         — application class + heterogeneous-tool reality
%%   S2  Research gap       — capability table + central RQ + sub-RQ keywords
%%   S3  Contribution       — what was built + forward pointers
%%%%%%%%%%%%%%%%%%%%%%%%%%%%%%%%%%%%%%%%%%%%%%%%%%%%%%%%%%%%%%%%%%%%%%%%%%%%%

\section{Introduction}

%---------------------------------------------------------------------------
\begin{frame}{Motivation}
\begin{columns}[T]
\begin{column}{0.58\textwidth}
\begin{itemize}
    \item Active mechatronic devices --- powered prostheses, exoskeletons --- are cyber-physical systems.
    \item Their performance depends on the coupling between device and the user's musculoskeletal system.
    \item Subsystem models live in different tools:
    \begin{itemize}
        \item mechanical structure: \ac{CAD} + \ac{FEM}
        \item control and electrical: Modelica
        \item musculoskeletal dynamics: OpenSim
    \end{itemize}
    \item Co-simulation couples these models; \ac{FMI} standardizes the simulation-unit interface.
\end{itemize}
\end{column}
\begin{column}{0.40\textwidth}
\centering
% TODO: insert a representative figure — e.g. device--human interaction
% sketch, or the case-study system overview from thesis/figures/6_case_study/
\textit{[figure placeholder]}
\end{column}
\end{columns}
\end{frame}

%---------------------------------------------------------------------------
\begin{frame}{Research Gap}
\begin{itemize}
    \item \ac{FMI}-based frameworks handle \ac{FMU} coupling and algebraic loops.
    \item None combines native \ac{FMU}, OpenSim, and \ac{FEM} backends with hybrid event handling and runtime model switching.
\end{itemize}

\vspace{0.3em}
\begin{table}[h]
\centering\footnotesize
\begin{tabular}{l c c c c c c}
\toprule
\textbf{Framework} & \textbf{\ac{FMU}} & \textbf{OpenSim} & \textbf{\ac{FEM}}
    & \textbf{Alg.\ loops} & \textbf{Hybrid} & \textbf{Switching} \\
\midrule
OMSimulator        & \checkmark & ---             & ---             & \checkmark & ($\checkmark$) & --- \\
INTO-CPS Maestro   & \checkmark & ---             & ---             & \checkmark & ($\checkmark$) & --- \\
CoFMPy             & \checkmark & ($\checkmark$)  & ($\checkmark$)  & \checkmark & ---             & --- \\
\textbf{\syssimx{}} & \checkmark & \checkmark     & \checkmark     & \checkmark & \checkmark     & \checkmark \\
\bottomrule
\end{tabular}
\end{table}

\vspace{0.4em}
{\small\textit{Central RQ:}\, How can a framework support heterogeneous hybrid co-simulation while providing consistent orchestration, event handling, and runtime switching between alternative models of the same subsystem?}

\vspace{0.2em}
{\footnotesize\textbf{RQ1} unified interface \,$\cdot$\, \textbf{RQ2} dependency analysis \,$\cdot$\, \textbf{RQ3} hybrid events \,$\cdot$\, \textbf{RQ4} model switching \,$\cdot$\, \textbf{RQ5} case-study evaluation}
\end{frame}

%---------------------------------------------------------------------------
\begin{frame}{Contribution}
\begin{block}{\syssimx{} --- a Python framework for heterogeneous hybrid co-simulation}
This thesis designs, implements, and evaluates:
\begin{itemize}
    \item a unified component abstraction for heterogeneous subsystem models,
    \item graph-based structural analysis (dependency detection, execution order, algebraic loops),
    \item hybrid event handling (detection, localization, dispatch, cascaded events),
    \item a multi-model component for runtime switching with state synchronization.
\end{itemize}
\end{block}

\begin{itemize}
    \item Evaluated on a controlled-pendulum case study combining \acp{FMU}, OpenSim, and \ac{FEM} models in one closed-loop scenario.
    \item Evidence path: architecture (Ch.~4) $\to$ implementation (Ch.~5) $\to$ case study (Ch.~6).
\end{itemize}
\end{frame}